\documentclass{article}
\usepackage{amsmath, amssymb, amsthm}
\usepackage{geometry}
\usepackage{physics}
\usepackage{graphicx}
\usepackage{booktabs}
\usepackage{hyperref}

\geometry{margin=1in}

\newtheorem{theorem}{Theorem}

\title{Projected Einasto Profile via Series Expansion\\
of Generalized Exponential Integrals}
\author{Johnny H.~Esteves}
\date{\today}

\begin{document}

\maketitle

\begin{abstract}
The Einasto density profile $\rho(r)=\rho_0\exp[-(r/h)^{1/n}]$ provides an
accurate description of dark-matter haloes in $N$-body simulations but has no
elementary closed form for its projected (line-of-sight) quantities.  We derive
exact, convergent series for the surface density $\Sigma(R)$, the projected
mass $M_{\rm 2D}(R)$, and the excess surface density $\Delta\Sigma(R)$,
expressed in terms of the generalized exponential integral $E_s(x)$ and the
lower incomplete gamma function $\gamma(s,x)$.  The series coefficients are
Catalan numbers scaled by $4^{-k}$, and the three quantities share a single
coefficient family through an integer ladder $\{k+1,1,k\}$.  We prove
convergence for all radii $R>0$ and shape indices $n>0$ via Raabe's test, with
algebraic term decay $O(k^{-3/2})$.
\end{abstract}

% ======================================================================
\section{Introduction}
\label{sec:intro}
% ======================================================================

The Einasto profile~\cite{Einasto1965,Navarro2004,Merritt2006} describes the
radial density of relaxed dark-matter haloes in cosmological $N$-body
simulations over several decades in radius.  The projected quantities---surface
density $\Sigma(R)$, cylindrical enclosed mass $M_{\rm 2D}(R)$, and excess
surface density $\Delta\Sigma(R)$---are central to gravitational lensing
analyses.  No elementary closed form exists for general shape index.

Retana-Montenegro et al.~\cite{RetanaMontenegro2012} expressed the surface
density in terms of the Fox~$H$ function via a Mellin integral transform, and
Dhar~\cite{Dhar2021} developed high-accuracy polynomial approximations using a
spherical-cap geometric decomposition.  In this work we show that the cap
decomposition yields \emph{exact} convergent series whose terms involve only
the generalized exponential integral $E_s(x)$---a standard special function
available in numerical libraries---with coefficients given by Catalan numbers.

% ======================================================================
\section{Definitions and Notation}
\label{sec:defs}
% ======================================================================

\subsection{The Einasto Profile}

We adopt the $(n,h)$ parametrization throughout:
\begin{equation}\label{eq:einasto}
\rho(r) = \rho_0 \exp\!\left[-\left(\frac{r}{h}\right)^{\!1/n}\right],
\qquad n>0,\; h>0\,.
\end{equation}
The mass enclosed in a sphere of radius $r$ is
\begin{equation}\label{eq:M3D}
M_{\rm 3D}(r) = 4\pi \rho_0\, n\, h^3\,
\gamma\!\left(3n,\;\left(\frac{r}{h}\right)^{\!1/n}\right),
\end{equation}
with total mass $M_{\rm tot} = 4\pi \rho_0\, n\, h^3\,\Gamma(3n)$.

\paragraph{Translation to $(\alpha, b, r_s)$.}
Setting $\alpha = 1/n$, $b = 2n$, and $h = r_s/(2n)^n$, one recovers
$\rho = \rho_0\exp[-b(r/r_s)^\alpha]$.

\subsection{Catalan Coefficients}
\label{sec:catalan}

The series coefficients are
\begin{equation}\label{eq:ck}
\boxed{\;
c_k \;=\; \frac{\mathrm{Cat}_k}{4^k}
    \;=\; \frac{1}{k+1}\binom{2k}{k}\frac{1}{4^k}
    \;=\; \frac{(2k-1)!!}{2^k\,(k+1)!}\,,
\qquad c_0=1\,.\;}
\end{equation}
They satisfy
\begin{equation}\label{eq:ck-properties}
\sum_{k=0}^{\infty} c_k = 2\,,
\qquad
c_k \sim \frac{1}{\sqrt{\pi}\,k^{3/2}} \quad (k\to\infty)\,.
\end{equation}
The Catalan generating function gives the key identity linking these
coefficients to the spherical-cap geometry (Sec.~\ref{sec:derivation}):
\begin{equation}\label{eq:catalan-gf}
1-\sqrt{1-u} \;=\; \frac{u}{2}\sum_{k=0}^{\infty} c_k\, u^k\,,
\qquad |u|\le 1\,.
\end{equation}

\subsection{Generalized Exponential Integral}

The generalized exponential integral is defined as
\begin{equation}\label{eq:Es-def}
E_s(x) = x^{s-1}\int_x^{\infty} \frac{e^{-t}}{t^s}\,\mathrm{d}t\,,
\qquad x>0\,,
\end{equation}
with the properties
\begin{align}
\frac{\mathrm{d}E_s}{\mathrm{d}x} &= -E_{s-1}(x)\,, \label{eq:Es-deriv}\\
s\,E_{s+1}(x) &= e^{-x} - x\,E_s(x)\,, \label{eq:Es-recurrence}
\end{align}
and the relation to the upper incomplete gamma function:
$E_s(x) = x^{s-1}\,\Gamma(1-s,x)$.

\subsection{Series Index}

Each term carries the exponential-integral index
\begin{equation}\label{eq:nuk}
\nu_k = 2kn - n + 1\,,
\end{equation}
which grows linearly with the summation index $k$.

% ======================================================================
\section{Main Results}
\label{sec:results}
% ======================================================================

\begin{theorem}[Projected Einasto profiles]\label{thm:main}
For the Einasto profile~\eqref{eq:einasto} with $n>0$ and $h>0$, the
projected quantities at cylindrical radius $R>0$ are:
\begin{align}
\Sigma(R) &= 2\rho_0\, n\, R \sum_{k=0}^{\infty}
(k+1)\,c_k\; E_{\nu_k}\!\left(\left(\frac{R}{h}\right)^{\!1/n}\right),
\label{eq:Sigma}\\[6pt]
M_{\rm 2D}(R) &= M_{\rm 3D}(R) + 2\pi\rho_0\, n\, R^3
\sum_{k=0}^{\infty} c_k\; E_{\nu_k}\!\left(\left(\frac{R}{h}\right)^{\!1/n}\right),
\label{eq:M2D}\\[6pt]
\Delta\Sigma(R) &= \frac{M_{\rm 3D}(R)}{\pi R^2}
- 2\rho_0\, n\, R \sum_{k=1}^{\infty}
k\,c_k\; E_{\nu_k}\!\left(\left(\frac{R}{h}\right)^{\!1/n}\right),
\label{eq:DeltaSigma}
\end{align}
where $c_k=\mathrm{Cat}_k/4^k$, $\nu_k=2kn-n+1$, and $M_{\rm 3D}$ is given
by~\eqref{eq:M3D}.
\end{theorem}

\noindent
The three quantities share the single coefficient family $c_k$ through an
\emph{integer ladder}:
\begin{center}
\begin{tabular}{lcc}
\toprule
Quantity & Weight & Origin \\
\midrule
$M_{\rm 2D}$ & $c_k$ & master (bare Catalan) \\
$\Sigma$ & $(k+1)\,c_k$ & radial derivative of $M_{\rm 2D}$ \\
$\Delta\Sigma$ & $k\,c_k$ & $\bar\Sigma - \Sigma$ \\
\bottomrule
\end{tabular}
\end{center}
The $k=0$ term drops out of $\Delta\Sigma$ automatically since
$k\,c_k\big|_{k=0}=0$.

% ======================================================================
\section{Derivation: Spherical Cap Decomposition}
\label{sec:derivation}
% ======================================================================

\subsection{Geometric Decomposition of $M_{\rm 2D}$}

Following Dhar~\cite{Dhar2021}, the projected mass decomposes as
\begin{equation}\label{eq:cap-split}
M_{\rm 2D}(R) = M_{\rm 3D}(R) + 2\int_R^{\infty}\rho(r)\,A(r)\,\mathrm{d}r\,,
\end{equation}
where
$A(r) = 2\pi r^2\bigl(1-\sqrt{1-R^2/r^2}\bigr)$
is the spherical-cap cross section at projected radius $R$ within a sphere of
radius $r$.

\subsection{Cap Area via Catalan Coefficients}

Using the generating-function identity~\eqref{eq:catalan-gf} with
$u=(R/r)^2$:
\begin{equation}
1-\sqrt{1-\frac{R^2}{r^2}}
= \frac{1}{2}\left(\frac{R}{r}\right)^{\!2}\sum_{k=0}^{\infty}c_k\left(\frac{R}{r}\right)^{\!2k},
\end{equation}
so the cap area simplifies to
\begin{equation}\label{eq:cap-area}
A(r) = \pi R^2 \sum_{k=0}^{\infty} c_k\left(\frac{R}{r}\right)^{\!2k}.
\end{equation}
The Catalan coefficients appear directly as the natural expansion of the
spherical-cap geometry.

\subsection{Integration Against the Einasto Profile}

Substituting~\eqref{eq:cap-area} into~\eqref{eq:cap-split}:
\begin{equation}\label{eq:DeltaM}
M_{\rm 2D}(R) - M_{\rm 3D}(R)
= 2\pi R^2 \sum_{k=0}^{\infty} c_k
\int_R^{\infty}\rho_0\,e^{-(r/h)^{1/n}}\left(\frac{R}{r}\right)^{\!2k}\mathrm{d}r\,.
\end{equation}
The key integral identity (proved in Appendix~\ref{app:integral}) is:
\begin{equation}\label{eq:Jk}
\boxed{\;
\int_R^{\infty}\rho_0\,e^{-(r/h)^{1/n}}\left(\frac{R}{r}\right)^{\!2k}\mathrm{d}r
= \rho_0\,n\,R\;E_{\nu_k}\!\left(\left(\frac{R}{h}\right)^{\!1/n}\right),
\qquad \nu_k = 2kn-n+1\,.\;}
\end{equation}
Substituting into~\eqref{eq:DeltaM} immediately gives
\begin{equation}
M_{\rm 2D}(R) = M_{\rm 3D}(R) + 2\pi\rho_0\,n\,R^3
\sum_{k=0}^{\infty}c_k\,E_{\nu_k}\!\left(\left(\frac{R}{h}\right)^{\!1/n}\right),
\end{equation}
establishing Eq.~\eqref{eq:M2D}.~$\square$

\subsection{Derivation of $\Sigma(R)$}

The surface density satisfies
$\Sigma(R) = (2\pi R)^{-1}\,\mathrm{d}M_{\rm 2D}/\mathrm{d}R$.
Differentiating term by term (justified by uniform convergence,
Sec.~\ref{sec:convergence}), using the chain rule and
recurrence~\eqref{eq:Es-recurrence}:
\begin{equation}
\frac{\mathrm{d}}{\mathrm{d}R}\left[R^3\,E_{\nu_k}\!\left(\left(\frac{R}{h}\right)^{\!1/n}\right)\right]
= 2(k+1)\,R^2\,E_{\nu_k}\!\left(\left(\frac{R}{h}\right)^{\!1/n}\right)
- \frac{R^2}{n}\,e^{-(R/h)^{1/n}},
\end{equation}
where we used $3+(\nu_k-1)/n = 3+(2k-1) = 2(k+1)$.  Summing over $k$ with
weight $c_k$ and multiplying by $\rho_0 n/R$:
\begin{equation}
\Sigma = \underbrace{2\rho_0 R\,e^{-(R/h)^{1/n}}}_{\text{from } M_{\rm 3D}}
+ 2\rho_0 nR\sum_{k=0}^{\infty}(k+1)c_k\,E_{\nu_k}
- \rho_0 R\,e^{-(R/h)^{1/n}}\underbrace{\sum_{k=0}^{\infty}c_k}_{=\,2}\,.
\end{equation}
The exponential terms cancel: $(2-2)\rho_0 R\,e^{-(R/h)^{1/n}}=0$, leaving
Eq.~\eqref{eq:Sigma}.~$\square$

\subsection{Derivation of $\Delta\Sigma(R)$}

The mean surface density is $\bar\Sigma(<\!R)=M_{\rm 2D}(R)/(\pi R^2)$.  Then
\begin{align}
\Delta\Sigma &= \bar\Sigma - \Sigma
= \frac{M_{\rm 3D}}{\pi R^2}
+ 2\rho_0 nR\sum_{k=0}^{\infty}c_k\,E_{\nu_k}
- 2\rho_0 nR\sum_{k=0}^{\infty}(k+1)c_k\,E_{\nu_k} \nonumber\\
&= \frac{M_{\rm 3D}}{\pi R^2}
- 2\rho_0 nR\sum_{k=0}^{\infty}k\,c_k\,E_{\nu_k}\!\left(\left(\frac{R}{h}\right)^{\!1/n}\right).
\end{align}
Since $k\,c_k|_{k=0}=0$, the sum starts at $k=1$, giving
Eq.~\eqref{eq:DeltaSigma}.~$\square$

% ======================================================================
\section{Convergence}
\label{sec:convergence}
% ======================================================================

We establish absolute convergence for all $R>0$ and $n>0$.  Write
$\Sigma = 2\rho_0 nR\sum_{k\ge0}u_k$ with
$u_k = (k+1)\,c_k\,E_{\nu_k}((R/h)^{1/n})$.

\paragraph{Coefficient asymptotics.}
From~\eqref{eq:ck-properties}, $(k+1)c_k\sim 1/(\sqrt\pi\,k^{1/2})$.

\paragraph{Exponential integral asymptotics.}
For large index at fixed argument $x=(R/h)^{1/n}>0$:
\begin{equation}
E_{\nu_k}(x) = \frac{e^{-x}}{x+\nu_k}\left[1+O(\nu_k^{-1})\right]
\;\xrightarrow{\;\nu_k\sim 2nk\;}\;
\frac{e^{-(R/h)^{1/n}}}{2nk}\,.
\end{equation}

\paragraph{Combined decay.}
\begin{equation}
u_k \;\sim\; \frac{e^{-(R/h)^{1/n}}}{2n\sqrt\pi}\;\frac{1}{k^{3/2}}
\qquad (k\to\infty).
\end{equation}
Convergence follows by comparison with $\sum k^{-3/2}$ ($p$-series, $p=3/2>1$).
The $M_{\rm 2D}$ terms decay as $k^{-5/2}$ and the $\Delta\Sigma$ terms as
$k^{-3/2}$---both convergent.

\paragraph{Raabe's test (sharp).}
The ordinary ratio test gives limit~1 (inconclusive, since decay is algebraic).
Raabe's test resolves:
\begin{equation}
\lim_{k\to\infty}\; k\left(\frac{u_k}{u_{k+1}}-1\right) = \frac{3}{2} > 1\,,
\end{equation}
establishing convergence and the exact algebraic order.  The prefactor
$e^{-(R/h)^{1/n}}$ ensures uniform convergence on bounded intervals and
accelerating convergence at large $R$.

% ======================================================================
\section{Special Cases}
\label{sec:special}
% ======================================================================

\subsection{Gaussian ($n=1/2$, $\alpha=2$)}

The profile becomes $\rho=\rho_0 e^{-(r/h)^2}$, with
$\nu_k = k+1/2$.  The series collapses to the known closed form:
\begin{equation}
\Sigma_{n=1/2}(R) = \sqrt\pi\,\rho_0\,h\;e^{-(R/h)^2}\,.
\end{equation}

\subsection{Exponential ($n=1$, $\alpha=1$)}

The profile becomes $\rho=\rho_0 e^{-r/h}$, with $\nu_k = 2k$.
The series recovers:
\begin{equation}
\Sigma_{n=1}(R) = 2\rho_0\,h\,(R/h)\,K_1(R/h)\,,
\end{equation}
where $K_1$ is the modified Bessel function of the second kind.

\subsection{First Three Terms}

\begin{center}
\begin{tabular}{ccccl}
\toprule
$k$ & $c_k$ & $(k+1)c_k$ & $\nu_k$ & Note \\
\midrule
0 & 1     & 1     & $1-n$  & $E_{1-n}(x)=x^{-n}\Gamma(n,x)$ \\
1 & $1/4$ & $1/2$ & $n+1$  & leading cap correction \\
2 & $1/8$ & $3/8$ & $3n+1$ & \\
\bottomrule
\end{tabular}
\end{center}
For typical dark-matter haloes ($n\simeq4$--$6$), the first 10--15 terms give
$10^{-6}$ relative accuracy at $R\sim h$.

% ======================================================================
\section{Numerical Implementation}
\label{sec:numerics}
% ======================================================================

Computing Eqs.~\eqref{eq:Sigma}--\eqref{eq:DeltaSigma} requires $E_s(x)$ for
real $s$:
\begin{itemize}
\item Integer $s\ge1$: \texttt{scipy.special.expn(s, x)}.
\item Real $s$: use $E_s(x)=x^{s-1}\Gamma(1-s,x)$ via
\texttt{mpmath.expint(s,x)} or the regularized incomplete gamma.
\item Truncation: from the $k^{-3/2}$ decay, $K$ terms give relative error
$\sim e^{(R/h)^{1/n}}/(2n\sqrt\pi\,K^{1/2})$.  At $R/h=1$, $n=4$: $K=15$
achieves $\epsilon<10^{-4}$.
\end{itemize}

% ======================================================================
\section{Conclusion}
\label{sec:conclusion}
% ======================================================================

We derived exact convergent series for the projected Einasto profile quantities
$\Sigma(R)$, $M_{\rm 2D}(R)$, and $\Delta\Sigma(R)$ in terms of the
generalized exponential integral $E_s(x)$.  The coefficients are Catalan
numbers $\mathrm{Cat}_k/4^k$, arising from the generating-function identity
$1-\sqrt{1-u}=(u/2)\sum c_k u^k$ that governs the spherical-cap geometry.
The three results share a single coefficient family through an integer ladder
$\{k+1,1,k\}$, with $M_{\rm 2D}$ as the master quantity.  Convergence is
established by Raabe's test with algebraic decay $O(k^{-3/2})$.

% ======================================================================
\appendix
\section{Integral Identity}
\label{app:integral}
% ======================================================================

We prove Eq.~\eqref{eq:Jk}.  Define
\begin{equation}
J_k = \int_R^{\infty}\rho_0\,e^{-(r/h)^{1/n}}\left(\frac{R}{r}\right)^{\!2k}\mathrm{d}r\,.
\end{equation}
Substitute $t=(r/h)^{1/n}$, so $r=h\,t^n$,
$\mathrm{d}r = n\,h\,t^{n-1}\,\mathrm{d}t$, and $(R/r)^{2k}=(R/h)^{2k}t^{-2nk}$:
\begin{align}
J_k &= \rho_0\,n\,h\left(\frac{R}{h}\right)^{\!2k}
\int_{(R/h)^{1/n}}^{\infty} e^{-t}\,t^{n-1-2nk}\,\mathrm{d}t \nonumber\\
&= \rho_0\,n\,h\left(\frac{R}{h}\right)^{\!2k}
\Gamma\!\left(n(1-2k),\;\left(\frac{R}{h}\right)^{\!1/n}\right).
\end{align}
Applying $\Gamma(a,x)=x^a\,E_{1-a}(x)$ with $a=n(1-2k)$ and $1-a=2nk-n+1=\nu_k$:
\begin{align}
J_k &= \rho_0\,n\,h\left(\frac{R}{h}\right)^{\!2k}
\left(\frac{R}{h}\right)^{\!(1-2k)}
E_{\nu_k}\!\left(\left(\frac{R}{h}\right)^{\!1/n}\right) \nonumber\\
&= \rho_0\,n\,R\;E_{\nu_k}\!\left(\left(\frac{R}{h}\right)^{\!1/n}\right),
\end{align}
where we used $x^{n(1-2k)}=((R/h)^{1/n})^{n(1-2k)}=(R/h)^{1-2k}$ and
$h(R/h)^{2k}(R/h)^{1-2k}=h(R/h)=R$.~$\square$

% ======================================================================
\section{Alternative Derivation via Abel Transform}
\label{app:abel}
% ======================================================================

The Abel transform gives $\Sigma$ directly:
\begin{equation}
\Sigma(R) = 2\int_R^{\infty}\rho(r)\,\frac{r}{\sqrt{r^2-R^2}}\,\mathrm{d}r
= 2\int_R^{\infty}\rho(r)\,\frac{1}{\sqrt{1-(R/r)^2}}\,\mathrm{d}r\,.
\end{equation}
Expanding $(1-u)^{-1/2}=\sum_{k\ge0}(k+1)c_k\,u^k$ with $u=(R/r)^2$ and
using the integral identity~\eqref{eq:Jk}:
\begin{equation}
\Sigma(R) = 2\sum_{k=0}^{\infty}(k+1)\,c_k\,J_k
= 2\rho_0\,n\,R\sum_{k=0}^{\infty}(k+1)\,c_k\,
E_{\nu_k}\!\left(\left(\frac{R}{h}\right)^{\!1/n}\right),
\end{equation}
recovering Eq.~\eqref{eq:Sigma} directly.  Both routes use the same integral
identity but enter through different binomial expansions:
\begin{itemize}
\item Cap route expands $1-\sqrt{1-u}=(u/2)\sum c_k u^k$ and integrates for
$M_{\rm 2D}$ (bare $c_k$), then differentiates to get $\Sigma$ (picks up
$k+1$).
\item Abel route expands $(1-u)^{-1/2}=\sum(k+1)c_k\,u^k$ and integrates
directly for $\Sigma$.
\end{itemize}
The identity $(k+1)c_k = \mathrm{d}/\mathrm{d}u[(u/2)\sum c_k u^k]|_{\text{coeff of }u^k}$ connects the two.

% ======================================================================
\section{Special Function Summary}
\label{app:special-fns}
% ======================================================================

\subsection{Incomplete Gamma Functions}
\begin{align}
\gamma(s,x) &= \int_0^x t^{s-1}e^{-t}\,\mathrm{d}t\,,\qquad
\Gamma(s,x) = \int_x^\infty t^{s-1}e^{-t}\,\mathrm{d}t\,,\\
\gamma(s,x) &+ \Gamma(s,x) = \Gamma(s) \qquad (s\neq0,-1,-2,\ldots)\,.
\end{align}

\subsection{Generalized Exponential Integral}
\begin{align}
E_s(x) &= x^{s-1}\int_x^\infty t^{-s}e^{-t}\,\mathrm{d}t
= x^{s-1}\,\Gamma(1-s,\,x)\,,\\
\frac{\mathrm{d}E_s}{\mathrm{d}x} &= -E_{s-1}(x)\,,\\
s\,E_{s+1}(x) &= e^{-x} - x\,E_s(x)\,.
\end{align}

\subsection{Catalan Generating Function}
\begin{equation}
\sum_{k=0}^{\infty}\mathrm{Cat}_k\,z^k = \frac{1-\sqrt{1-4z}}{2z}\,,
\quad |z|\le\tfrac14\,.
\end{equation}
Setting $z=u/4$: $\;\sum c_k u^k = 2(1-\sqrt{1-u})/u$, hence
$1-\sqrt{1-u}=(u/2)\sum c_k u^k$ and $\sum c_k=2$ (at $u=1$).

% ======================================================================
\begin{thebibliography}{99}

\bibitem{Einasto1965}
J.~Einasto, Trudy Astrofiz. Inst. Alma-Ata \textbf{5}, 87 (1965).

\bibitem{Navarro2004}
J.~F.~Navarro et al., Mon. Not. R. Astron. Soc. \textbf{349}, 1039 (2004).

\bibitem{Merritt2006}
D.~Merritt, A.~W.~Graham, B.~Moore, J.~Diemand, and B.~Terzi\'c,
Astron. J. \textbf{132}, 2685 (2006).

\bibitem{RetanaMontenegro2012}
E.~Retana-Montenegro, E.~Van Hese, G.~Gentile, M.~Baes, and
F.~Frutos-Alfaro, Astron. Astrophys. \textbf{540}, A70 (2012).

\bibitem{Dhar2021}
B.~K.~Dhar, Mon. Not. R. Astron. Soc. \textbf{504}, 4583 (2021).

\end{thebibliography}

\end{document}
